\documentclass[11pt,a4paper]{article}
\usepackage[T1]{fontenc}
\usepackage{lmodern,array,amsmath,amssymb,booktabs,enumitem,fancyhdr,lastpage,needspace,microtype}
\usepackage[margin=19mm,headheight=14pt,footskip=11mm]{geometry}
\usepackage[hidelinks]{hyperref}
\setlength{\parindent}{0pt}
\setlength{\parskip}{5pt}
\setlength{\emergencystretch}{2em}
\pagestyle{fancy}\fancyhf{}
\fancyhead[L]{\small CS182 --- Lecture 04}
\fancyhead[R]{\small Personal answer copy}
\fancyfoot[L]{\small Independent practice quiz}
\fancyfoot[R]{\small Page \thepage\ of \pageref{LastPage}}
\newcommand{\SetAnswer}[2]{\expandafter\gdef\csname QuizAnswer#1\endcsname{#2}}
\newcommand{\GetAnswer}[1]{\csname QuizAnswer#1\endcsname}
\newlength{\BlankAnswerHeight}\setlength{\BlankAnswerHeight}{16mm}
\newcommand{\AnswerBlank}{\leavevmode\par\vspace{\BlankAnswerHeight}}
% 本文件为独立可填写答题卷；用 pdfLaTeX 编译两次。
% 在 BEGIN TF ANSWERS 和 BEGIN MC ANSWERS 中填写短答案。
% 大题请搜索 BEGIN ANSWER Q25(a) 等注释；每个小问都有独立区域。
% 写完小问后删除该处的 \AnswerBlank，或将 BlankAnswerHeight 改为 0mm。
% 所有打印内容及你填写的回答请使用英文。中文只出现在源文件注释中。
% build_per_lecture.py 不会覆盖已经存在的 Quiz.tex，以保留你的作答。
% BEGIN STUDENT INFO / 个人信息
\newcommand{\StudentName}{}
\newcommand{\StudentID}{}
% END STUDENT INFO
% BEGIN TF ANSWERS / 判断题答题区（第二组花括号填 T 或 F）
\SetAnswer{TF01}{} % Question 1
\SetAnswer{TF02}{} % Question 2
\SetAnswer{TF03}{} % Question 3
\SetAnswer{TF04}{} % Question 4
\SetAnswer{TF05}{} % Question 5
\SetAnswer{TF06}{} % Question 6
\SetAnswer{TF07}{} % Question 7
\SetAnswer{TF08}{} % Question 8
\SetAnswer{TF09}{} % Question 9
\SetAnswer{TF10}{} % Question 10
\SetAnswer{TF11}{} % Question 11
\SetAnswer{TF12}{} % Question 12
% END TF ANSWERS
% BEGIN MC ANSWERS / 不定项选择答题区（填选项字母，如 BD；仅格式示例）
\SetAnswer{MC13}{} % Question 13
\SetAnswer{MC14}{} % Question 14
\SetAnswer{MC15}{} % Question 15
\SetAnswer{MC16}{} % Question 16
\SetAnswer{MC17}{} % Question 17
\SetAnswer{MC18}{} % Question 18
\SetAnswer{MC19}{} % Question 19
\SetAnswer{MC20}{} % Question 20
\SetAnswer{MC21}{} % Question 21
\SetAnswer{MC22}{} % Question 22
\SetAnswer{MC23}{} % Question 23
\SetAnswer{MC24}{} % Question 24
% END MC ANSWERS

\begin{document}
\begin{center}{\Large\bfseries CS182: Introduction to Machine Learning}\par{\Large\bfseries Lecture 04: Optimization Primer}\par{\large Lecture Quiz and Answer Copy}\end{center}
\textbf{Scope:} Optimization\_Primer.pdf. Independent practice material based on the official 2026 Fall lecture.\par\textbf{Total:} 100 marks. \quad \textbf{Suggested time:} 120 minutes.
\par\textbf{Name:} \underline{\makebox[60mm][l]{\StudentName}}\hfill\textbf{Student ID:} \underline{\makebox[45mm][l]{\StudentID}}
\begin{center}\begin{tabular}{lrrr}\toprule
Section & Questions & Each & Marks\\\midrule
A. Misconceptions: True/False & 12 & 1 & 12\\
B. Multiple response & 12 & 2 & 24\\
C. Proofs and formula derivations & 4 & 10 & 40\\
D. Applications & 2 & 12 & 24\\\midrule
Total & 30 & & 100\\\bottomrule\end{tabular}\end{center}
\textbf{Scoring.} In Section B each item has exactly four options, of which one,
two, or three are correct. The complete correct set earns 2 marks; a nonempty
proper subset earns 1 mark; any incorrect option or no selection earns 0.
T/F answers earn 1 mark only when correct. Show reasoning in Sections C and D;
part marks are available for valid intermediate work. Equivalent correct methods receive credit.
\par\medskip\textbf{Instructions.} Answer in English. Use natural logarithms.
Exact values and justified approximations are accepted. All vectors are columns
unless stated otherwise. Use local assumptions, label conventions, and tie rules.
\par\medskip\textbf{Answer sheet --- Section A}
\begin{center}\renewcommand{\arraystretch}{1.15}\begin{tabular}{|c|c|c|c|c|c|}
\hline
1 & 2 & 3 & 4 & 5 & 6\\\hline
\makebox[20mm]{\rule{0pt}{5mm}\GetAnswer{TF01}} & \makebox[20mm]{\rule{0pt}{5mm}\GetAnswer{TF02}} & \makebox[20mm]{\rule{0pt}{5mm}\GetAnswer{TF03}} & \makebox[20mm]{\rule{0pt}{5mm}\GetAnswer{TF04}} & \makebox[20mm]{\rule{0pt}{5mm}\GetAnswer{TF05}} & \makebox[20mm]{\rule{0pt}{5mm}\GetAnswer{TF06}}\\\hline
\end{tabular}\end{center}
\begin{center}\renewcommand{\arraystretch}{1.15}\begin{tabular}{|c|c|c|c|c|c|}
\hline
7 & 8 & 9 & 10 & 11 & 12\\\hline
\makebox[20mm]{\rule{0pt}{5mm}\GetAnswer{TF07}} & \makebox[20mm]{\rule{0pt}{5mm}\GetAnswer{TF08}} & \makebox[20mm]{\rule{0pt}{5mm}\GetAnswer{TF09}} & \makebox[20mm]{\rule{0pt}{5mm}\GetAnswer{TF10}} & \makebox[20mm]{\rule{0pt}{5mm}\GetAnswer{TF11}} & \makebox[20mm]{\rule{0pt}{5mm}\GetAnswer{TF12}}\\\hline
\end{tabular}\end{center}
\par\medskip\textbf{Answer sheet --- Section B}
\begin{center}\renewcommand{\arraystretch}{1.15}\begin{tabular}{|c|c|c|c|c|c|}
\hline
13 & 14 & 15 & 16 & 17 & 18\\\hline
\makebox[20mm]{\rule{0pt}{5mm}\GetAnswer{MC13}} & \makebox[20mm]{\rule{0pt}{5mm}\GetAnswer{MC14}} & \makebox[20mm]{\rule{0pt}{5mm}\GetAnswer{MC15}} & \makebox[20mm]{\rule{0pt}{5mm}\GetAnswer{MC16}} & \makebox[20mm]{\rule{0pt}{5mm}\GetAnswer{MC17}} & \makebox[20mm]{\rule{0pt}{5mm}\GetAnswer{MC18}}\\\hline
\end{tabular}\end{center}
\begin{center}\renewcommand{\arraystretch}{1.15}\begin{tabular}{|c|c|c|c|c|c|}
\hline
19 & 20 & 21 & 22 & 23 & 24\\\hline
\makebox[20mm]{\rule{0pt}{5mm}\GetAnswer{MC19}} & \makebox[20mm]{\rule{0pt}{5mm}\GetAnswer{MC20}} & \makebox[20mm]{\rule{0pt}{5mm}\GetAnswer{MC21}} & \makebox[20mm]{\rule{0pt}{5mm}\GetAnswer{MC22}} & \makebox[20mm]{\rule{0pt}{5mm}\GetAnswer{MC23}} & \makebox[20mm]{\rule{0pt}{5mm}\GetAnswer{MC24}}\\\hline
\end{tabular}\end{center}
\newpage\section*{Section A. True/False \hfill 12 marks}
\par\noindent\begin{minipage}{\linewidth}
\textbf{1.} The feasible set defined by convex inequalities \(f_i(x)\le0\) and affine equalities is convex.
\end{minipage}\par\vspace{3mm}
\par\noindent\begin{minipage}{\linewidth}
\textbf{2.} For a symmetric matrix \(Q\), the unconstrained quadratic \(f(x)=\frac12x^TQx+c^Tx\) is convex if and only if \(Q\succeq0\).
\end{minipage}\par\vspace{3mm}
\par\noindent\begin{minipage}{\linewidth}
\textbf{3.} The Lagrange dual function is concave even when the primal objective or inequality constraints are nonconvex.
\end{minipage}\par\vspace{3mm}
\par\noindent\begin{minipage}{\linewidth}
\textbf{4.} If a constraint is active at a KKT point, its multiplier must be strictly positive.
\end{minipage}\par\vspace{3mm}
\par\noindent\begin{minipage}{\linewidth}
\textbf{5.} If a primal feasible point and a dual feasible point have the same finite objective value, they certify primal and dual optimality.
\end{minipage}\par\vspace{3mm}
\par\noindent\begin{minipage}{\linewidth}
\textbf{6.} A differentiable function is convex whenever its gradient vanishes at at least one point.
\end{minipage}\par\vspace{3mm}
\par\noindent\begin{minipage}{\linewidth}
\textbf{7.} Every local minimizer of a convex function over a convex feasible set is a global minimizer.
\end{minipage}\par\vspace{3mm}
\par\noindent\begin{minipage}{\linewidth}
\textbf{8.} For differentiable convex objectives, every positive constant gradient-descent step size guarantees convergence from every starting point.
\end{minipage}\par\vspace{3mm}
\par\noindent\begin{minipage}{\linewidth}
\textbf{9.} Equality multipliers may have either sign because the corresponding constraint values are zero at every primal feasible point.
\end{minipage}\par\vspace{3mm}
\par\noindent\begin{minipage}{\linewidth}
\textbf{10.} Every differentiable convex problem with an attained optimum has an optimal finite multiplier vector.
\end{minipage}\par\vspace{3mm}
\par\noindent\begin{minipage}{\linewidth}
\textbf{11.} At a KKT point, a strictly positive inequality multiplier implies that the corresponding inequality is active.
\end{minipage}\par\vspace{3mm}
\par\noindent\begin{minipage}{\linewidth}
\textbf{12.} For a minimization problem with constraints \(f_i(x)\le0\), every real multiplier vector provides a valid dual lower bound.
\end{minipage}\par\vspace{3mm}
\newpage\section*{Section B. Multiple Response \hfill 24 marks}
\textit{Select all correct options. Each question has 1--3 correct options.}\par
\par\noindent\begin{minipage}{\linewidth}
% BEGIN QUESTION Q13 / 题目与选项区
\textbf{13.} Let \(f\) be differentiable and convex on \(\mathbb R^n\). Select every conclusion that follows.
\begin{enumerate}[label=(\Alph*),leftmargin=8mm,itemsep=2pt,topsep=2pt]
\item The Hessian must be positive definite wherever it exists.
\item If \(\nabla f(x)=0\), then \(x\) is a global unconstrained minimizer.
\item \(f(y)\ge f(x)+\nabla f(x)^T(y-x)\) for all \(x,y\).
\item Every global minimizer is unique.
\end{enumerate}
% END QUESTION Q13

\end{minipage}\par\vspace{3mm}
\par\noindent\begin{minipage}{\linewidth}
% BEGIN QUESTION Q14 / 题目与选项区
\textbf{14.} For a general primal minimization problem, select every statement justified without assuming convexity or strong duality.
\begin{enumerate}[label=(\Alph*),leftmargin=8mm,itemsep=2pt,topsep=2pt]
\item The dual function is concave on its effective domain.
\item A zero gradient at a primal feasible point guarantees the global constrained optimum.
\item Any primal feasible \(x\) and dual feasible multipliers satisfy \(g(\lambda,\nu)\le f_0(x)\).
\item If a differentiable Lagrangian has zero gradient at \(x\), that point necessarily realizes its infimum.
\end{enumerate}
% END QUESTION Q14

\end{minipage}\par\vspace{3mm}
\par\noindent\begin{minipage}{\linewidth}
% BEGIN QUESTION Q15 / 题目与选项区
\textbf{15.} Consider minimizing \((x-3)^2\) subject to \(x\le1\), with \(L(x,\lambda)=(x-3)^2+\lambda(x-1)\). Select every correct calculation.
\begin{enumerate}[label=(\Alph*),leftmargin=8mm,itemsep=2pt,topsep=2pt]
\item \(g(\lambda)=2\lambda-\lambda^2/4\).
\item The unconstrained minimizer of \(L(\cdot,\lambda)\) is \(x=3-\lambda/2\).
\item The dual optimum is attained at \(\lambda=4\), with value \(4\).
\item \(\lambda=-2\) gives a valid dual certificate because the Lagrangian infimum is finite.
\end{enumerate}
% END QUESTION Q15

\end{minipage}\par\vspace{3mm}
\par\noindent\begin{minipage}{\linewidth}
% BEGIN QUESTION Q16 / 题目与选项区
\textbf{16.} Which formulations satisfy the definition of a convex optimization problem? All variables shown are real.
\begin{enumerate}[label=(\Alph*),leftmargin=8mm,itemsep=2pt,topsep=2pt]
\item Minimize \(-x^2\) subject to \(-1\le x\le1\).
\item Minimize \(x\) subject to \(x^2=1\).
\item Minimize \(x^2+y^2\) subject to \(x+y=1\) and \(x\ge0\).
\item Minimize \(3x-y\) subject to \(x^2+y^2\le4\).
\end{enumerate}
% END QUESTION Q16

\end{minipage}\par\vspace{3mm}
\newpage\section*{Section B. Multiple Response \hfill 24 marks}
\textit{Select all correct options. Each question has 1--3 correct options.}\par
\par\noindent\begin{minipage}{\linewidth}
% BEGIN QUESTION Q17 / 题目与选项区
\textbf{17.} Let \(G=\begin{pmatrix}2&6\\0&4\end{pmatrix}\) and \(f(x)=\tfrac12x^TGx\) on \(\mathbb R^2\). Select every correct statement.
\begin{enumerate}[label=(\Alph*),leftmargin=8mm,itemsep=2pt,topsep=2pt]
\item \(\nabla f(x)=(2x_1+3x_2,\,3x_1+4x_2)^T\).
\item \(f\) is convex because the diagonal entries of \(G\) are positive.
\item The Hessian equals the nonsymmetric matrix \(G\).
\item The stationary point \(x=0\) is a saddle point.
\end{enumerate}
% END QUESTION Q17

\end{minipage}\par\vspace{3mm}
\par\noindent\begin{minipage}{\linewidth}
% BEGIN QUESTION Q18 / 题目与选项区
\textbf{18.} Select every correct statement about stopping gradient descent or checking optimality.
\begin{enumerate}[label=(\Alph*),leftmargin=8mm,itemsep=2pt,topsep=2pt]
\item The norm \(\|\nabla f(x)\|\) can be used as a scalar stationarity residual.
\item A small gradient norm alone gives a universal numerical bound on objective error for every differentiable convex function.
\item A zero gradient certifies a global unconstrained minimizer for a differentiable convex objective.
\item For a constrained minimizer on the boundary, the objective gradient must itself be zero.
\end{enumerate}
% END QUESTION Q18

\end{minipage}\par\vspace{3mm}
\par\noindent\begin{minipage}{\linewidth}
% BEGIN QUESTION Q19 / 题目与选项区
\textbf{19.} For minimizing \(f_0(x)\) subject to \(f_i(x)\le0\) and \(h_j(x)=0\), select every correct Lagrangian or dual convention.
\begin{enumerate}[label=(\Alph*),leftmargin=8mm,itemsep=2pt,topsep=2pt]
\item The dual function is obtained by minimizing the Lagrangian only over primal feasible points.
\item The dual problem maximizes \(g(\lambda,\nu)\) over \(\lambda\ge0\).
\item If the infimum of the Lagrangian is \(-\infty\), that multiplier vector is dual feasible under the lecture definition.
\item \(L=f_0+\sum_i\lambda_if_i+\sum_j\nu_jh_j\), with \(\lambda\ge0\) and unrestricted \(\nu\).
\end{enumerate}
% END QUESTION Q19

\end{minipage}\par\vspace{3mm}
\par\noindent\begin{minipage}{\linewidth}
% BEGIN QUESTION Q20 / 题目与选项区
\textbf{20.} The primal linear program is \(\min_x c^Tx\) subject to \(Ax\le b\), where \(A\in\mathbb R^{m\times n}\) and \(x\) is unrestricted. Select every correct statement.
\begin{enumerate}[label=(\Alph*),leftmargin=8mm,itemsep=2pt,topsep=2pt]
\item The dual has \(n\) variables and \(m\) equality constraints.
\item Its dual maximizes \(-b^T\lambda\) subject to \(A^T\lambda+c=0\) and \(\lambda\ge0\).
\item If \(A^T\lambda+c\ne0\), the dual function equals \(-\infty\).
\item The nonnegativity restriction on \(\lambda\) can be dropped because \(x\) is unrestricted.
\end{enumerate}
% END QUESTION Q20

\end{minipage}\par\vspace{3mm}
\newpage\section*{Section B. Multiple Response \hfill 24 marks}
\textit{Select all correct options. Each question has 1--3 correct options.}\par
\par\noindent\begin{minipage}{\linewidth}
% BEGIN QUESTION Q21 / 题目与选项区
\textbf{21.} For \(f(x)=\tfrac12(x_1^2+4x_2^2)\), gradient descent uses a constant step size \(\eta=0.6\). Select every correct conclusion.
\begin{enumerate}[label=(\Alph*),leftmargin=8mm,itemsep=2pt,topsep=2pt]
\item The iterates converge to zero from every starting point because the objective is strictly convex.
\item The update is \(x_{k+1}=(0.4x_{1,k},-1.4x_{2,k})^T\).
\item At any starting point with \(x_{2,0}\ne0\), the second coordinate changes sign each iteration and grows in magnitude.
\item If \(x_{2,0}=0\), the iterates converge to zero.
\end{enumerate}
% END QUESTION Q21

\end{minipage}\par\vspace{3mm}
\par\noindent\begin{minipage}{\linewidth}
% BEGIN QUESTION Q22 / 题目与选项区
\textbf{22.} Let \(f\) be convex and let \(X\) be a finitely supported random vector in its convex domain. Select every valid use of convexity.
\begin{enumerate}[label=(\Alph*),leftmargin=8mm,itemsep=2pt,topsep=2pt]
\item The pointwise maximum of two convex functions is convex.
\item The pointwise minimum of any two convex functions is convex.
\item Jensen's inequality requires independent coordinates of \(X\).
\item \(f(E[X])\ge E[f(X)]\) for every such convex function and random vector.
\end{enumerate}
% END QUESTION Q22

\end{minipage}\par\vspace{3mm}
\par\noindent\begin{minipage}{\linewidth}
% BEGIN QUESTION Q23 / 题目与选项区
\textbf{23.} Minimize \(x^2\) subject to \(-x\le0\) and \(-2x\le0\). Select every correct statement at \(x^*=0\).
\begin{enumerate}[label=(\Alph*),leftmargin=8mm,itemsep=2pt,topsep=2pt]
\item \((x^*,\lambda_1,\lambda_2)=(0,1,0)\) satisfies all KKT conditions.
\item Both inequalities are active.
\item The stationarity equation is \(2x-\lambda_1-\lambda_2=0\).
\item Both multipliers must be strictly positive.
\end{enumerate}
% END QUESTION Q23

\end{minipage}\par\vspace{3mm}
\par\noindent\begin{minipage}{\linewidth}
% BEGIN QUESTION Q24 / 题目与选项区
\textbf{24.} An algorithm has a primal feasible point with objective \(11.2\) and dual feasible multipliers with dual value \(10.9\). Select every justified conclusion.
\begin{enumerate}[label=(\Alph*),leftmargin=8mm,itemsep=2pt,topsep=2pt]
\item The optimal value must equal the midpoint \(11.05\).
\item The distance from the primal point to every minimizer is at most \(0.3\).
\item The primal point is at most \(0.3\) above the optimal objective value.
\item The optimal primal value belongs to \([10.9,11.2]\).
\end{enumerate}
% END QUESTION Q24

\end{minipage}\par\vspace{3mm}
\newpage\section*{Section C. Proofs and Formula Derivations}
\textbf{25. Proving the first-order characterization of convexity}\hfill\textbf{[10 marks]}\par\medskip
\begin{enumerate}[label=(\alph*),leftmargin=8mm,itemsep=3pt]
\item Let \(f\) be differentiable on an open convex domain. Starting with the defining convexity inequality, prove that convexity implies \(f(y)\ge f(x)+\nabla f(x)^T(y-x)\) for any \(x,y\) in the domain.\hfill\textbf{[5 marks]}
\item Conversely, assume the first-order inequality holds for every pair \(x,y\). Prove the defining convexity inequality using \(z=(1-t)x+ty\), \(0<t<1\), as the expansion point.\hfill\textbf{[5 marks]}
\end{enumerate}\par\medskip\hrule\medskip\textbf{Student response}\par
% BEGIN ANSWER Q25 / 第25题答题区
\begin{enumerate}[label=(\alph*),leftmargin=8mm,itemsep=3pt]
\item
% >>> BEGIN ANSWER Q25(a)
% 在这里写英文推导或证明；完成后可删除下方空白占位命令。
\AnswerBlank
% <<< END ANSWER Q25(a)
\item
% >>> BEGIN ANSWER Q25(b)
% 在这里写英文推导或证明；完成后可删除下方空白占位命令。
\AnswerBlank
% <<< END ANSWER Q25(b)
\end{enumerate}
% END ANSWER Q25
\newpage\section*{Section C. Proofs and Formula Derivations}
\textbf{26. Dual concavity, weak duality, and complementary slackness}\hfill\textbf{[10 marks]}\par\medskip
\begin{enumerate}[label=(\alph*),leftmargin=8mm,itemsep=3pt]
\item For \(L(x,\lambda,\nu)=f_0(x)+\sum_i\lambda_if_i(x)+\sum_j\nu_jh_j(x)\), prove directly that \(g(\lambda,\nu)=\inf_x L(x,\lambda,\nu)\) is concave, without imposing convexity on the primal functions. Work with multiplier pairs having finite dual values.\hfill\textbf{[4 marks]}
\item For primal feasible \(x\) and \(\lambda\ge0\), derive \(g(\lambda,\nu)\le f_0(x)\) and hence weak duality. Explain the different restrictions on inequality and equality multipliers.\hfill\textbf{[3 marks]}
\item Suppose feasible \(x^*,\lambda^*,\nu^*\) have finite values and \(f_0(x^*)=g(\lambda^*,\nu^*)\). Derive complementary slackness term by term and show that \(x^*\) minimizes the Lagrangian.\hfill\textbf{[3 marks]}
\end{enumerate}\par\medskip\hrule\medskip\textbf{Student response}\par
% BEGIN ANSWER Q26 / 第26题答题区
\begin{enumerate}[label=(\alph*),leftmargin=8mm,itemsep=3pt]
\item
% >>> BEGIN ANSWER Q26(a)
% 在这里写英文推导或证明；完成后可删除下方空白占位命令。
\AnswerBlank
% <<< END ANSWER Q26(a)
\item
% >>> BEGIN ANSWER Q26(b)
% 在这里写英文推导或证明；完成后可删除下方空白占位命令。
\AnswerBlank
% <<< END ANSWER Q26(b)
\item
% >>> BEGIN ANSWER Q26(c)
% 在这里写英文推导或证明；完成后可删除下方空白占位命令。
\AnswerBlank
% <<< END ANSWER Q26(c)
\end{enumerate}
% END ANSWER Q26
\newpage\section*{Section C. Proofs and Formula Derivations}
\textbf{27. Why convex KKT conditions are sufficient}\hfill\textbf{[10 marks]}\par\medskip
\begin{enumerate}[label=(\alph*),leftmargin=8mm,itemsep=3pt]
\item Assume \(f_0,f_i\) are differentiable convex functions on \(\mathbb R^n\), and \(h_j\) are affine. Suppose \((x^*,\lambda^*,\nu^*)\) satisfies primal feasibility, \(\lambda^*\ge0\), complementary slackness, and stationarity. Prove that \(x^*\) globally minimizes \(L(\cdot,\lambda^*,\nu^*)\).\hfill\textbf{[4 marks]}
\item Use this Lagrangian minimum to prove primal and dual optimality and zero duality gap. State whether a constraint qualification is needed for this sufficiency implication once the displayed KKT point is given.\hfill\textbf{[3 marks]}
\item If \(f_0\) is strictly convex, prove that the primal minimizer is unique. Identify the role of the convexity of the feasible set.\hfill\textbf{[3 marks]}
\end{enumerate}\par\medskip\hrule\medskip\textbf{Student response}\par
% BEGIN ANSWER Q27 / 第27题答题区
\begin{enumerate}[label=(\alph*),leftmargin=8mm,itemsep=3pt]
\item
% >>> BEGIN ANSWER Q27(a)
% 在这里写英文推导或证明；完成后可删除下方空白占位命令。
\AnswerBlank
% <<< END ANSWER Q27(a)
\item
% >>> BEGIN ANSWER Q27(b)
% 在这里写英文推导或证明；完成后可删除下方空白占位命令。
\AnswerBlank
% <<< END ANSWER Q27(b)
\item
% >>> BEGIN ANSWER Q27(c)
% 在这里写英文推导或证明；完成后可删除下方空白占位命令。
\AnswerBlank
% <<< END ANSWER Q27(c)
\end{enumerate}
% END ANSWER Q27
\newpage\section*{Section C. Proofs and Formula Derivations}
\textbf{28. Deriving a quadratic gradient-descent step-size bound}\hfill\textbf{[10 marks]}\par\medskip
\begin{enumerate}[label=(\alph*),leftmargin=8mm,itemsep=3pt]
\item Let \(f(x)=\frac12x^TQx+c^Tx+r\), where \(Q\) is symmetric. Derive its gradient and Hessian by expanding \(f(x+h)\). State the condition for convexity on \(\mathbb R^n\).\hfill\textbf{[3 marks]}
\item Assume \(0\preceq Q\preceq LI\) with \(L>0\), let \(g=\nabla f(x)\), and define \(x^+=x-\eta g\). Derive an upper bound on \(f(x^+)-f(x)\), and give a range of \(\eta\) guaranteeing strict decrease whenever \(g\ne0\).\hfill\textbf{[4 marks]}
\item Now assume \(Q\succ0\). Prove convergence of constant-step gradient descent for \(0<\eta<2/\lambda_{\max}(Q)\) by diagonalizing \(Q\). Give a one-dimensional counterexample at the upper endpoint.\hfill\textbf{[3 marks]}
\end{enumerate}\par\medskip\hrule\medskip\textbf{Student response}\par
% BEGIN ANSWER Q28 / 第28题答题区
\begin{enumerate}[label=(\alph*),leftmargin=8mm,itemsep=3pt]
\item
% >>> BEGIN ANSWER Q28(a)
% 在这里写英文推导或证明；完成后可删除下方空白占位命令。
\AnswerBlank
% <<< END ANSWER Q28(a)
\item
% >>> BEGIN ANSWER Q28(b)
% 在这里写英文推导或证明；完成后可删除下方空白占位命令。
\AnswerBlank
% <<< END ANSWER Q28(b)
\item
% >>> BEGIN ANSWER Q28(c)
% 在这里写英文推导或证明；完成后可删除下方空白占位命令。
\AnswerBlank
% <<< END ANSWER Q28(c)
\end{enumerate}
% END ANSWER Q28
\newpage\section*{Section D. Applications}
\textbf{29. Capacity planning with an active resource cap}\hfill\textbf{[12 marks]}\par\medskip
\begin{enumerate}[label=(\alph*),leftmargin=8mm,itemsep=3pt]
\item A planning model chooses nonnegative production amounts \(x,y\) with total \(x+y=4\), subject to the capacity cap \(x\le2.5\). It minimizes \(f(x,y)=(x-5)^2+2(y-2)^2\). Eliminate the equality constraint, determine the feasible interval, and find the unique minimizer and objective value.\hfill\textbf{[4 marks]}
\item Write the Lagrangian using multipliers \(\alpha,\beta,\lambda\ge0\) for \(-x,-y,x-2.5\le0\), respectively, and unrestricted \(\nu\) for \(x+y-4=0\). Find multipliers at the proposed optimum and give a dual certificate by evaluating the Lagrangian infimum.\hfill\textbf{[4 marks]}
\item Before convergence, a solver reports the feasible plan \((x,y)=(2,2)\) and multipliers \((\alpha,\beta,\lambda,\nu)=(0,0,2,2)\). Compute their dual value, a guaranteed interval for \(p^*\), and an upper bound on the plan's objective error.\hfill\textbf{[4 marks]}
\end{enumerate}\par\medskip\hrule\medskip\textbf{Student response}\par
% BEGIN ANSWER Q29 / 第29题答题区
\begin{enumerate}[label=(\alph*),leftmargin=8mm,itemsep=3pt]
\item
% >>> BEGIN ANSWER Q29(a)
% 在这里写英文推导或证明；完成后可删除下方空白占位命令。
\AnswerBlank
% <<< END ANSWER Q29(a)
\item
% >>> BEGIN ANSWER Q29(b)
% 在这里写英文推导或证明；完成后可删除下方空白占位命令。
\AnswerBlank
% <<< END ANSWER Q29(b)
\item
% >>> BEGIN ANSWER Q29(c)
% 在这里写英文推导或证明；完成后可删除下方空白占位命令。
\AnswerBlank
% <<< END ANSWER Q29(c)
\end{enumerate}
% END ANSWER Q29
\newpage\section*{Section D. Applications}
\textbf{30. Choosing a stable learning rate for a quadratic calibration loss}\hfill\textbf{[12 marks]}\par\medskip
\begin{enumerate}[label=(\alph*),leftmargin=8mm,itemsep=3pt]
\item A calibration model minimizes \(f(u,v)=\tfrac12(2u^2+8v^2)-4u+8v\) on \(\mathbb R^2\). Compute the gradient, Hessian, unique optimizer, and minimum value. Explain why a stationary point is sufficient here.\hfill\textbf{[4 marks]}
\item From \((u_0,v_0)=(0,0)\), calculate one gradient-descent step and its objective for \(\eta=0.2\) and for \(\eta=0.3\). Compare each objective to the initial value and to the optimum.\hfill\textbf{[4 marks]}
\item Derive the constant-step interval guaranteeing convergence from every start. At the boundary step \(\eta=0.25\), determine the behavior from \((0,0)\) and explain why a zero change in objective would not certify convergence.\hfill\textbf{[4 marks]}
\end{enumerate}\par\medskip\hrule\medskip\textbf{Student response}\par
% BEGIN ANSWER Q30 / 第30题答题区
\begin{enumerate}[label=(\alph*),leftmargin=8mm,itemsep=3pt]
\item
% >>> BEGIN ANSWER Q30(a)
% 在这里写英文推导或证明；完成后可删除下方空白占位命令。
\AnswerBlank
% <<< END ANSWER Q30(a)
\item
% >>> BEGIN ANSWER Q30(b)
% 在这里写英文推导或证明；完成后可删除下方空白占位命令。
\AnswerBlank
% <<< END ANSWER Q30(b)
\item
% >>> BEGIN ANSWER Q30(c)
% 在这里写英文推导或证明；完成后可删除下方空白占位命令。
\AnswerBlank
% <<< END ANSWER Q30(c)
\end{enumerate}
% END ANSWER Q30
\end{document}
